\section{Implementation And Evaluation}\label{sec:impl}

We have implemented a tool based on the above approach, called
\name{}, that targets distributed system that modeled as actors (i.e., message-passing over components). \name{} consists of approximately 14K lines of OCaml code and uses Z3~\cite{de2008z3} as its backend solver for SMT queries.\footnote{Our \techreport{} provides a docker image that contains the source code of \name{} and all our benchmarks.}

\paragraph{Evaluation setting} \name{} takes a safety property in s-LTL$_f$ that is expected to be violated, and handler specification $\Delta$ that is a set of messages labeled with generatable or observable as well as corresponding \PAT{}s (e.g., the types shown in \autoref{fig:spec}). The typing context used by \name{} includes signatures for a number of arithmetic primitives, including the pure operators listed in
\autoref{fig:term-syntax}.
During synthesis, \name{} will first negate the given property as an violation property (e.g., $\negSCA$ in \autoref{sec:overview}), and consecutively explore controllers that guides the components that consistent with $\Delta$ to produce traces consistent with the violation property. Each explored controller indicates a specific message order, and \name{} returns union of these controller as the final synthesized result.

The synthesized controller is directly executable with a interpreter implementing the operational semantics in \autoref{fig:selected-semantics}, where the auxiliary handler semantics ($\alpha \vDash \zevent{op}{\overline{c}} \Downarrow \beta$) can be implemented as a random sampler to provide messages $\beta$ that align with handler specifications. However, to evaluate our approach in a more realistic setting, we actually adapted our synthesized controller to run alongside a third-party message handler implementation.
Specifically, we used handlers implemented in the P language\cite{?}, a trace-based, actor-model programming language tailored for modeling distributed systems and testing user-defined safety properties. Similar to our setup, P language monitors inter-component messages, dynamically producing a trace which it evaluates against safety criteria. Unlike our method, however, P’s runtime scheduler explores arbitrary message interleavings randomly during execution.
To test our synthesized controller with P language handlers while retaining scheduling control, we compile the controller as a special component in P that works like a ``coordinator''. In this setup, messages between components are first routed to our controller, where they are buffered and then being forwarded to the actual handlers. The forwarding order is determined by the $\mykw{obs}$ statements in the controller program, allowing it to control the message scheduling.

Our experimental evaluation addresses three research questions:
\begin{itemize}
\item[\textbf{Q1}:] Is \name{} \emph{expressive}? Can it synthesizes controller for a diverse range of protocols in distributed system with variety of safety properties?
\item[\textbf{Q2}:] Is \name{} \emph{effective}? Do synthesized controller programs realize violation of safety property and are also in a reasonable size?
\item[\textbf{Q3}:] Is \name{} \emph{efficient}? Is it able to synthesize controller programs in a reasonable amount of time?
\end{itemize} %

\begin{table}[ht!]
  \renewcommand{\arraystretch}{0.8}
  % \vspace*{-.1in}
  \caption{\small Experimental results of using \name{} to synthesize controller for reactive distributed systems. Benchmarks from prior works are annotated with their source: P Language~\cite{?}($^{\dagger}$), ModP~\cite{?}($^{\diamond}$), and MessageChain~\cite{?}($^{\star}$). We also include the an real world model (\textsf{Kermit2PCModel}$^{\square}$) written in P language that models 2PC protocol used in Amazon Web Services (AWS). The components under test are original written in P language, where the handler specification are manually defined as \PAT{}s. The synthesized controller program is then compiled into P language and run with components under test in P language runtime. \name{} can synthesize union of controllers that indicates one specific message order. For efficiency result, we set a $2$ min time bound where the synthesis time (i.e., t$_\text{total}$) is the average time to find a single controller. }
\vspace*{-.1in}
\footnotesize
\setlength{\tabcolsep}{2.3pt}
\begin{tabular}{c||cc||ccc|c||c|ccccc}
  \toprule
 Benchmark & \#$\eff{op}$ & \#qualifier & \#var &(\#$\mykw{gen}$, \#$\mykw{obs}$) & \#$\mykw{assert}$ & \#try & t$_\text{total}$(s) & t$_\text{SAT}$(s) & t$_\text{refine}$(s) & \#SAT & \#FW & \#BW \\
 \midrule
\textsf{Database} & $4$ & $7$ & $8$ & ($3$, $3$) & $4$ & $2.0$ & $6.02$ & $5.99$ & $5.88$ & $905$ & $4$ & $6$\\
\midrule
\textsf{EspressoMachine}$^{\dagger}$ & $13$ & $3$ & $1$ & ($2$, $8$) & $1$ & $4.0$ & $1.15$ & $1.12$ & $0.99$ & $165$ & $2$ & $11$\\
\midrule
\textsf{Simplified2PC}$^{\dagger}$ & $9$ & $14$ & $7$ & ($2$, $6$) & $5$ & $1.9$ & $6.83$ & $6.80$ & $6.66$ & $1043$ & $3$ & $8$\\
\midrule
\textsf{HeartBeat}$^{\dagger}$ & $7$ & $18$ & $9$ & ($4$, $10$) & $9$ & $1.0$ & $7.05$ & $6.99$ & $6.46$ & $1073$ & $4$ & $20$\\
\midrule
\textsf{BankServer}$^{\dagger}$ & $6$ & $17$ & $15$ & ($2$, $3$) & $5$ & $1.0$ & $8.24$ & $8.21$ & $7.08$ & $1191$ & $2$ & $5$\\
\midrule
\textsf{RingLeaderElection}$^{\star}$ & $3$ & $19$ & $12$ & ($2$, $6$) & $8$ & $1.0$ & $6.37$ & $6.32$ & $6.18$ & $964$ & $2$ & $18$\\
\midrule
\textsf{Firewall}$^{\star}$ & $5$ & $17$ & $12$ & ($2$, $8$) & $9$ & $1.0$ & $5.26$ & $5.21$ & $5.07$ & $788$ & $5$ & $8$\\
\midrule
\textsf{ChainReplication}$^{\diamond}$ & $7$ & $32$ & $26$ & ($4$, $9$) & $10$ & $1.0$ & $27.06$ & $26.89$ & $26.76$ & $4016$ & $6$ & $19$\\
\midrule
\textsf{Paxos}$^{\diamond}$ & $10$ & $30$ & $36$ & ($4$, $10$) & $13$ & $1.0$ & $59.71$ & $59.27$ & $59.24$ & $8763$ & $4$ & $16$\\
\midrule
\textsf{Raft} & $9$ & $26$ & $29$ & ($3$, $14$) & $14$ & $1.0$ & $55.62$ & $55.14$ & $55.25$ & $8356$ & $10$ & $22$\\
\midrule
\textsf{Kermit2PCModel}$^{\square}$ & $17$ & $70$ & $36$ & ($4$, $10$) & $10$ & $1.0$ & $63.69$ & $63.33$ & $63.46$ & $9023$ & $6$ & $12$\\
\bottomrule
\end{tabular}

\label{tab:evaluation}
% \vspace*{-.15in}
\end{table}

We have evaluated \name{} on a corpus of complex reactive distributed system model drawn from a variety of sources (shown in the title of \autoref{tab:evaluation}), including molding of variety of applications such as database (the motivating example in \autoref{sec:overview}), espresso machine, banker server, and firewall, as well as well-known distributed protocols and algorithm including two-phase commit (2PC)\cite{?}, ring leader election\cite{?}, chain replication\cite{?}, Paxos\cite{?}, and Raft\cite{?}(\textbf{Q1}). We expect these reactive distributed system model to violate variety safety properties, include Read-Your-Writes (WRY) policy introduced in \autoref{sec:overview}, strong consistency, and unique leader invariants. We would like to highlight that \PAT{} is able to express communication over multiple nodes that playing the similar roles. For example, the \textsc{RingLeaderElection} models the election algorithm elects a leader among a set of nodes. We add additional source ($\I{src}$) and destination ($\I{dest}$) into the message payload to distinguish messages from different nodes. Then, handler of message $\eff{elect}$ has the following \PAT{}:
{\small\begin{align*}
    &\Code{n}\hasoty{tNode}{\top} \sarr \Code{m}\hasoty{tNode}{\top} \sarr \Code{ld}\hasoty{tNode}{v = \Code{m}} \sarr
    \\&\rg{\globalA\evparenth{\top}}{\lastA\msgB{elect}{\I{src}\;\I{dest}\;\I{leader}}{ \I{src} = \Code{n} \land \I{dest} = \Code{m} \land \I{leader} = \Code{m} }}{\finalA\msgB{beLeader}{\I{leader}}{\I{leader} = \Code{m}}}
\end{align*}}\noindent
which capture the following principle in ring leader election: when a node $\Code{m}$ received a message that elects ($\eff{elect}$) itself as leader ($\I{leader} = \Code{m}$) from another node $\Code{n}$, then node $\Code{m}$ can announce itself to be the leader ($\eff{beLeader}$).
Details about the reactive distributed system and safety properties used in our benchmarks are provided in \techreport{}.

\autoref{tab:evaluation} divides the results of our experiments into
three three categories, separated by the double bars. The first gives the complexity of benchmarks with respect to the number of messages (\#$\eff{op}$) involved and the number of qualifiers used in the both \PAT{}s and safety property that expressed in s-LTL$_f$. Our approach is able to model these reactive distributed system with $3$ - $17$ different messages and the handler specification sand safety property uses $3$ - $70$ qualifiers (\textbf{Q1}). The second group of columns describes the size of synthesized controller, including the number of variables (\#var), generative messages (\#$\mykw{gen}$), observable messages (\#$\mykw{obs}$), and assertions (\#$\mykw{assert}$) used in the controller program. Our synthesized controllers has $1$ - $36$ variables, $5$ - $17$ messages in total, and $1$ - $14$ assertions, which is a reasonable size. Note that, the size of synthesized program is roughly proportional to the complexity of the benchmarks (\textbf{Q2}), where the number of qualifiers most interfere the number of variables and assertions, while the controller with more effect applications caused by the total number of message in the system. As mentioned in \autoref{fig:algo}, our algorithm bias to the shorter controller program and avoiding the events that has no contribution to the safety property violation. Thus, the controller uses the messages in an efficient way where the total effect applications is close to total number of message in the system. The last column (\#try) in the second group reports the average number of try times that a synthesized program to violate the safety property in dynamic. For all reactive system that has deterministic handlers, the synthesized controller doesn't need to try (\#try is equal to $1$). On the other hand, for example, the handler of $\eff{writeReq}$ in \textsc{Database} introduced in \autoref{sec:overview} is allowed to accept of reject the write request, which is modeled as an non-deterministic choice with equal distribution in the implementation written in P language, which leads $2$ runs to violate the RYW policy.

The last group of columns in \autoref{tab:evaluation} reports the complexity of our synthesis algorithm. The first column indicate total synthesis time (t$_\text{total}$), which shows that \name{} performs quite reasonably on all of our benchmarks, taking between $1.15$ to $63.69$ seconds for synthesis a controller (\textbf{Q3}). As expected, more complicated reactive distributed system (i.e., those with larger values in the \#$\eff{op}$ and \#qualifier columns), take longer to synthesis. The last two columns indicate time cost by SAT queries solving (t$_\text{SAT}$) and the property refinement loop (t$_\text{refine}$), which has close number to the total synthesis time. It means that the most of time cost comes from the SAT queries, moreover, the property refinement loop take much more time than term derivation, which is consistent with our previous discussion in \autoref{sec:algo}. The last three columns additionally analyze the behaviors of \name{} with respect to the number of SAT queries (\#SAT), as well as forward synthesis (\#FW) and backward synthesis step (\#BW) preformed by the property refinement loop. Unsurprisingly, generating more
SMT queries (\#SAT) result in a longer synthesis time, where \#SAT is depends on the iteration of property refinement loop (i.e., the sum of \#FW and \#BW). Intuitively, the synthesizing longer program (the sum of \#$\mykw{gen}$ and \#$\mykw{obs}$) needs more iterations. Often, the forward and backward synthesis steps exceed the number of messages in the controller program in most of benchmarks, as \name{} may need to backtrack when a wrong type or message interleaving is selected, which future iterations cannot resolve.

\paragraph{Case study} We aim to demonstrate that \name{} can be applied to real-world scenarios, as illustrated in the \textsf{Kermit2PCModel} benchmark, which models the two-phase commit (2PC) protocol used in Amazon S3 storage systems \cite{?} in the P language. The original model checks a consistency property for transactions in 2PC, specifically that if there exists a key $\Code{k}$ updated within a transaction $\Code{i}$, any successful read response to the user should return the latest written value $\Code{x}$ from that transaction. In our approach, the corresponding violation property is defined as follows:
{\small\begin{align*}
    \finalA (&\msg{updateRsp}{tid = \Code{i} \land key = \Code{k} \land v = \Code{x} \land st = \Code{OK}} \land 
    \\&\nextA \neg \msg{updateRsp}{tid = \Code{i} \land key = \Code{k} \land st = \Code{OK}} \untilA \msg{readRsp}{tid = \Code{i} \land key = \Code{k} \land v \not= \Code{x} \land st = \Code{OK}})
\end{align*}}\noindent
where the field $tid$ represents the transaction id, while other fields have the same meanings as in the example from \autoref{sec:intro}. This scenario is challenging because the violation requires observing a successful $\eff{readRsp}$ with an incorrect value ($v \neq \Code{x}$). This misread can only occur when the preceding $\eff{readReq}$ selects a valid transaction id $\Code{i}$, which is the storage system manages rather than the user. Without explicitly requesting a new transaction, random testing has virtually no chance of ``guessing'' a valid transaction id. However, since the \PAT for $\eff{readReq}$ includes a history automaton $\finalA \msg{startTxnRsp}{tid = \Code{i}}$, requiring that the user previously received a valid transaction id $\Code{i}$, \name{} can synthesize a controller program that strategically requests a new transaction to trigger the intended violation.